% META: {"id": "TEXP-ARI-EX-047", "chapitre": "TEXP-ARITHMETIQUE", "type_objet": "exercice", "capacites": ["TEXP-ARI-C2"], "capacites_codes": ["C2"], "methodes": ["M3"], "parcours": 2, "competences": ["raisonner"], "duree_min": 20, "mode_creation": "ex_nihilo", "sources_inspiration": [], "fichier_tex": "chapitres/TEXP-ARITHMETIQUE/exercices/TEXP-ARI-EX-047.tex", "corrige_tex": "chapitres/TEXP-ARITHMETIQUE/corriges/TEXP-ARI-CO-047.tex", "status": "approved"}
% BEGIN-VERIFY
% from sympy import *
% x = symbols('x')
% # f(x) = ln(x), concave => f(x) <= f(a) + f'(a)(x-a) pour tout a
% # en a=1: f(1)=0, f'(1)=1, tangente y=x-1 => ln(x) <= x-1
% # application: ln(1.1) <= 0.1
% import math
% assert math.log(1.1) < 0.1
% assert math.log(1.1) > 0.05  # approximation x/2
% # f(x) = -x^2 convexe? Non, concave. f''=-2 < 0
% # En a=0: tangente y=0, -x^2 <= 0 toujours vrai
% g = -x**2 - 0  # f - tangente
% assert expand(g) == -x**2
% assert g.subs(x, 0) == 0
% assert g.subs(x, 1) == -1
% END-VERIFY
\begin{exercice}{TEXP-ARI-EX-047}{2}{20}
En utilisant la concavite de $\ln$, majorer $\ln(1{,}1)$ par une valeur simple, puis minorer $\ln(1{,}1)$ a l'aide d'une autre tangente bien choisie.

\end{exercice}
